\documentclass[11pt,letterpaper]{report}

\usepackage[utf8]{inputenc}
\usepackage[T1]{fontenc}
\usepackage[spanish,provide=*]{babel}
\usepackage{amsmath,amssymb}
\usepackage{booktabs}
\usepackage{array}
\usepackage{graphicx}
\usepackage{float}
\usepackage{caption}
\usepackage[margin=2.5cm]{geometry}
\usepackage{enumitem}
\usepackage[dvipsnames]{xcolor}

\renewcommand{\arraystretch}{1.15}

\begin{document}

\begin{center}
{\LARGE\bfseries\linespread{1.3}\selectfont 2. GENERACIÓN DE NÚMEROS RECTANGULARES\par}
\end{center}
\vspace{0.5cm}

\setcounter{chapter}{2}
\setcounter{page}{19}

%==========================================================
En todos los experimentos de simulación existe la necesidad de generar
valores de variables aleatorias que representan a una cierta
distribución de probabilidad. Durante un experimento de simulación, el
proceso de generar un valor de la variable aleatoria de una distribución
particular, puede repetirse tantas veces como se desee y tantas veces
como distribuciones de probabilidad existan en el experimento de
simulación. Sin embargo, es conveniente señalar que el proceso de
generación de variables aleatorias no uniformes se hace a partir de la
generación de números rectangulares. Por consiguiente, el objetivo de
este capítulo es mostrar un panorama general de las diferentes técnicas
que existen para generar números rectangulares.

La importancia de los números rectangulares (distribución uniforme)
radica en su uso para la generación de variables aleatorias más
complicadas que son requeridas en los experimentos de simulación.
Algunos autores como Tocher, han sugerido tres formas para obtener los
números rectangulares: La provisión externa, la generación interna a
partir de un proceso físico al azar y la generación interna de
sucesiones de dígitos por medio de una relación de recurrencia. El
primer método implica tener los números aleatorios, como por ejemplo las
tablas de la Rand, en una cinta magnética o en un disco y tratar a estos
números como datos de entrada para el problema que se está simulando.
El segundo método implica utilizar algún aditamento especial de la
computadora digital capaz de registrar los resultados de un proceso
aleatorio y además, reduzca esos resultados a sucesiones de dígitos. El
tercer método, y uno de los más aceptados, implica la generación de
estos números rectangulares a través de una relación de recurrencia.

Independientemente del proceso o procedimiento que se utilice para la
generación de los números rectangulares, estos deben de poseer ciertas
características deseables que aseguren o aumenten la confiabilidad de
los resultados obtenidos de la simulación. Tales características son:

\begin{enumerate}
\item Uniformemente distribuidos.
\item Estadísticamente independientes.
\item Reproducibles.
\item Período largo (sin repetición dentro de una longitud determinada
de la sucesión).
\item Generados a través de un método rápido.
\item Generados a través de un método que no requiera mucha capacidad
de almacenamiento de la computadora.
\end{enumerate}

Finalmente, es necesario señalar que algunos autores califican a los
números rectangulares generados a través de relaciones de recurrencia
con números \emph{pseudoaleatorios}, por ser una sucesión de dígitos
generada mediante una regla puramente determinística. Sin embargo, esta
objeción puede superarse, al menos parcialmente, al tomar el punto de
vista un tanto pragmático de que una sucesión puede considerarse
aleatoria si satisface un cierto conjunto de pruebas estadísticas de
aleatoriedad.

%==========================================================
\section*{2.1. Generadores congruenciales lineales}

Varios esquemas han sido propuestos para la generación de los números
pseudoaleatorios a través de relaciones matemáticas de recurrencia.
Estos números se consideran pseudoaleatorios, porque aunque pasan todas
las pruebas estadísticas de aleatoriedad, ellos son de hecho
\textcolor{BurntOrange}{completamente determinísticos}. Actualmente, casi todas las computadoras
incluyen en sus programas de biblioteca alguna variante de los métodos
congruenciales sugeridos por Lehmer. Los dos métodos congruenciales más
populares son: congruencial mixto y congruencial multiplicativo.

\subsection*{2.1.1. Congruencial mixto}

Los generadores congruenciales lineales generan una secuencia de números
pseudoaleatorios en la cual el próximo número pseudoaleatorio es
determinado a partir del último número generado, es decir, el número
pseudoaleatorio $X_{n+1}$ es derivado a partir del número pseudoaleatorio
$X_n$. Para el caso particular del generador congruencial mixto, la
relación de recurrencia es la siguiente:

\begin{equation}
X_{n+1} = (aX_n + c) \bmod m
\end{equation}

donde:
\begin{align*}
X_0 &= \text{la semilla } (X_0 > 0) \\
a &= \text{el multiplicador } (a > 0) \\
c &= \text{constante aditiva } (c > 0) \\
m &= \text{el módulo } (m > X_0,\ m > a\ \text{y}\ m > c)
\end{align*}

Esta relación de recurrencia nos dice que $X_{n+1}$ es el residuo de
dividir $aX_n + c$ entre el módulo. Lo anterior significa que los
valores posibles de $X_{n+1}$ son $0, 1, 2, 3, \ldots, m-1$, es decir,
$m$ representa el número posible de valores diferentes que pueden ser
generados.

Con el propósito de ilustrar la generación de números pseudoaleatorios a
través de este método, suponga que se tiene un generador en el cual los
valores de sus parámetros son: $a=5$, $c=7$, $X_0=4$ y $m=8$. Para estos
valores, la secuencia de números pseudoaleatorios y números uniformes
$(X_{n+1}/m)$ son mostrados en la tabla 2.1. Como se puede apreciar en
esta tabla, el período del generador es 8.

Después de haber analizado este ejemplo, podría pensarse que el período
de todo generador es siempre igual a $m$. Sin embargo, esto es falso
porque el período depende de los valores asignados a los parámetros $a$,
$c$, $X_0$ y $m$, es decir, se requiere seleccionar valores adecuados
para estos parámetros con el fin de que el generador tenga período
completo.

\begin{table}[H]
\centering
\caption*{\textbf{Tabla 2.1} Números pseudoaleatorios del generador $X_{n+1} = (5X_n + 7) \bmod 8$}
\begin{tabular}{ccccc}
\toprule
$n$ & $X_n$ & $(5X_n+7)/8$ & $X_{n+1}$ & Números uniformes \\
\midrule
0 & 4 & $3 + 3/8$ & 3 & $3/8$ \\
1 & 3 & $2 + 6/8$ & 6 & $6/8$ \\
2 & 6 & $4 + 5/8$ & 5 & $5/8$ \\
3 & 5 & $4 + 0/8$ & 0 & $0$ \\
4 & 0 & $0 + 7/8$ & 7 & $7/8$ \\
5 & 7 & $5 + 2/8$ & 2 & $2/8$ \\
6 & 2 & $2 + 1/8$ & 1 & $1/8$ \\
7 & 1 & $1 + 4/8$ & 4 & $4/8$ \\
\bottomrule
\end{tabular}
\end{table}

Para ilustrar el caso que se presenta cuando el período $< m$, suponga
que se tiene un generador en el cual los valores de sus parámetros son:
$a = X_0 = c = 7$ y $m = 10$. Para estos valores, la secuencia de
números pseudoaleatorios y números uniformes son mostrados en la tabla
2.2. Como puede apreciarse en esta tabla, el período del generador es 4.
Esto demuestra que una selección inadecuada de los valores de los
parámetros del generador, puede conducirnos a obtener resultados
indeseables y poco confiables del experimento de simulación.

\begin{table}[H]
\centering
\caption*{\textbf{Tabla 2.2} Números pseudoaleatorios del generador $X_{n+1} = (7X_n + 7) \bmod 10$}
\begin{tabular}{ccccc}
\toprule
$n$ & $X_n$ & $(7X_n+7)/10$ & $X_{n+1}$ & Números uniformes \\
\midrule
0 & 7 & $5 + 6/10$ & 6 & $6/10$ \\
1 & 6 & $4 + 9/10$ & 9 & $9/10$ \\
2 & 9 & $7 + 0/10$ & 0 & $0$ \\
3 & 0 & $0 + 7/10$ & 7 & $7/10$ \\
\bottomrule
\end{tabular}
\end{table}

De los ejemplos anteriores, se advierte la necesidad de establecer
algunas reglas que puedan ser utilizadas en la selección de los valores
de los parámetros, para que el generador resultante tenga período
completo. Algunas de estas reglas se mencionan a continuación:

\subsubsection*{a) Selección de $m$}

Existen dos opciones para seleccionar el valor apropiado del módulo:

\begin{enumerate}
\item Seleccionar $m$ de modo que sea el número primo más grande
posible y a la vez sea menor que $p^d$, donde $p$ es la base del sistema
(binario, decimal, hexadecimal, etc.) que se está utilizando y $d$ es el
número de bits que tiene una palabra de computadora en ese sistema. Por
ejemplo, si se tiene una computadora IBM 370 que trabaja en sistema
binario, entonces $p=2$ y $d=32$.
\item Seleccionar $m$ como $p^d$. Cuando $m$ toma este valor se facilita
el cálculo del número rectangular ($U_n = X_n/m$), ya que sólo se corre
el punto binario o decimal a la izquierda del número. Sin embargo, se ha
comprobado que cuando el módulo toma este valor, los últimos dígitos del
número pseudoaleatorio generado no se comportan en forma aleatoria.
\end{enumerate}

Para ilustrar el problema que se presenta cuando se utiliza el criterio
2, suponga que se tiene un generador cuyos parámetros son: $a=81$,
$c=89$, $X_0=5$ y $m=10^2$. Para estos valores, la secuencia de números
pseudoaleatorios son mostrados en la tabla 2.3. En esta tabla se puede
apreciar que el último dígito del número pseudoaleatorio tiene un
período de 10. Esto significa que el último dígito puede ser determinado
a partir de la siguiente relación de recurrencia:

\begin{equation}
Y_{n+1} = (Y_n + 9) \bmod 10
\end{equation}

\begin{table}[H]
\centering
\caption*{\textbf{Tabla 2.3} Números pseudoaleatorios del generador $X_{n+1} = (81X_n + 89) \bmod 100$}
\begin{tabular}{cccccccccc}
\toprule
$n$ & $X_n$ & $n$ & $X_n$ & $n$ & $X_n$ & $n$ & $X_n$ & $n$ & $X_n$ \\
\midrule
1 & 94 & 21 & 74 & 41 & 54 & 61 & 34 & 81 & 14 \\
2 & 03 & 22 & 83 & 42 & 63 & 62 & 43 & 82 & 23 \\
3 & 32 & 23 & 12 & 43 & 92 & 63 & 72 & 83 & 52 \\
4 & 81 & 24 & 61 & 44 & 41 & 64 & 21 & 84 & 01 \\
5 & 50 & 25 & 30 & 45 & 10 & 65 & 90 & 85 & 70 \\
6 & 39 & 26 & 19 & 46 & 99 & 66 & 79 & 86 & 59 \\
7 & 48 & 27 & 28 & 47 & 08 & 67 & 88 & 87 & 68 \\
8 & 77 & 28 & 57 & 48 & 37 & 68 & 17 & 88 & 97 \\
9 & 26 & 29 & 06 & 49 & 86 & 69 & 66 & 89 & 46 \\
10 & 95 & 30 & 75 & 50 & 55 & 70 & 35 & 90 & 15 \\
11 & 84 & 31 & 64 & 51 & 44 & 71 & 24 & 91 & 04 \\
12 & 93 & 32 & 73 & 52 & 53 & 72 & 33 & 92 & 13 \\
13 & 22 & 33 & 02 & 53 & 82 & 73 & 62 & 93 & 42 \\
14 & 71 & 34 & 51 & 54 & 31 & 74 & 11 & 94 & 91 \\
15 & 40 & 35 & 20 & 55 & 00 & 75 & 80 & 95 & 60 \\
16 & 29 & 36 & 09 & 56 & 89 & 76 & 69 & 96 & 49 \\
17 & 38 & 37 & 18 & 57 & 98 & 77 & 78 & 97 & 58 \\
18 & 67 & 38 & 47 & 58 & 27 & 78 & 07 & 98 & 87 \\
19 & 16 & 39 & 96 & 59 & 76 & 79 & 56 & 99 & 36 \\
20 & 85 & 40 & 65 & 60 & 45 & 80 & 25 & 100 & 05 \\
\bottomrule
\end{tabular}
\end{table}

Del ejemplo anterior, es posible generalizar una relación de recurrencia
que relacione los últimos dígitos del número pseudoaleatorio generado.
Si $m=p^d$, se ha encontrado que la relación de recurrencia de los
últimos dígitos es la siguiente:

\begin{equation}
Y_{n+1,i} = X_{n+1} \bmod p^i, \qquad i < d
\end{equation}

donde:
\begin{align*}
Y_{n+1,i} &= \text{últimos } i \text{ dígitos del número pseudoaleatorio } X_{n+1}. \\
i &= \text{últimos } i \text{ dígitos que se están considerando. El valor de } i \text{ puede ser } 1, 2, 3, \ldots, d-1.
\end{align*}

Por ejemplo, si $i=1$, la expresión anterior permite determinar el valor
del último dígito del número pseudoaleatorio $X_{n+1}$, si $i=2$, se
determina el valor de los dos últimos dígitos y así sucesivamente.

La expresión anterior puede ser manipulada y expresada en función de los
últimos dígitos del número pseudoaleatorio inmediato anterior, de
acuerdo a lo siguiente:

\[
Y_{n+1,i} = X_{n+1} \bmod p^i
\]

sustituyendo en la expresión anterior la relación de recurrencia
$X_{n+1} = (aX_n+c) \bmod m$ se obtiene:

\[
Y_{n+1,i} = (aX_n + c) \bmod p^d \bmod p^i
\]

pero si $q$ es la parte entera que resulta de dividir $(aX_n+c)/p^d$,
entonces la expresión anterior se transforma en:

\begin{align*}
Y_{n+1,i} &= (aX_n + c - qp^d) \bmod p^i \\
Y_{n+1,i} &= (aX_n + c) \bmod p^i \\
Y_{n+1,i} &= aX_n \bmod p^i + c \bmod p^i
\end{align*}

y como $X \bmod m = (X \bmod m) \bmod m$, entonces la expresión anterior
se transforma en:

\[
Y_{n+1,i} = (aX_n \bmod p^i) \bmod p^i + c \bmod p^i
\]

y haciendo uso de la identidad anterior, la expresión anterior se
transforma finalmente en:

\begin{equation}
Y_{n+1,i} = (aY_{n,i} + c) \bmod p^i
\end{equation}

esta última expresión también tendrá período completo ($p^i$) si el
generador tiene un período de $p^d$.

\subsubsection*{b) Selección de $a$}

El valor seleccionado de $a$ debe ser entero impar, y además no debe ser
divisible por 3 ó 5. Sin embargo, si queremos asegurar que el generador
tenga período completo, el valor de $a$ se debe de seleccionar de
acuerdo al siguiente criterio:

\begin{align*}
(a-1) \bmod 4 &= 0 \quad \text{si } 4 \text{ es un factor de } m \\
(a-1) \bmod b &= 0 \quad \text{si } b \text{ es un factor primo de } m
\end{align*}

Usualmente se selecciona $a$ como $2^k+1$ cuando se trabaja en sistema
binario y $10^k+1$ cuando se trabaja en sistema decimal. En ambos casos
el valor de $k$ debe ser mayor o igual a 2.

\subsubsection*{c) Selección de $c$}

El valor seleccionado para este parámetro puede ser cualquier constante.
Sin embargo, si se desean asegurar buenos resultados el valor de $c$
debe ser $c \bmod 8 = 5$ si se trabaja en sistema binario y $c \bmod
200 = 21$ si se trabaja en sistema decimal. Más específicamente, el
valor de $c$ debe ser un entero impar y relativamente primo a $m$.

\subsubsection*{d) Selección de $X_0$}

Para el generador congruencial mixto, se ha encontrado que el valor de
la semilla es irrelevante, es decir, el valor de este parámetro resulta
tener poca o ninguna influencia sobre las propiedades estadísticas de
las sucesiones.

Finalmente, antes de terminar la discusión de este generador, conviene
señalar que existen otras formas matemáticas de representarlo. Tales
formas son las siguientes:

\begin{equation}
X_n = \left\{ a^n X_0 + c \left\{ \frac{a^n - 1}{a-1} \right\} \right\} \bmod m
\end{equation}

\begin{equation}
X_{n+k} = \left\{ a^n X_k + c \left\{ \frac{a^n - 1}{a-1} \right\} \right\} \bmod m
\end{equation}

Con la primera expresión el $n$-ésimo número pseudoaleatorio se obtiene
a partir de la semilla. Con la segunda expresión el $n+k$-ésimo número
pseudoaleatorio se obtiene a partir del $k$-ésimo número, es decir, si
por ejemplo $n+k=10$ y $k=4$, entonces significa que el número
pseudoaleatorio 10 se va a obtener a partir del número 4.

\subsection*{2.1.2. Congruencial multiplicativo}

Al igual que el generador congruencial mixto, el generador congruencial
multiplicativo determina el próximo número pseudoaleatorio a partir del
último número generado, de acuerdo a la siguiente relación de
recurrencia:

\begin{equation}
X_{n+1} = aX_n \bmod m
\end{equation}

Para este generador se recomienda también seleccionar adecuadamente los
valores de los parámetros $a$, $X_0$ y $m$, con el fin de asegurar un
período máximo para las sucesiones generadas por este método. Los
valores de estos parámetros dependerán del sistema en que se trabaje, es
decir, estos parámetros tomarán valores distintos si se trabaja en
sistema decimal, que si se trabaja en sistema binario. Por consiguiente,
a continuación se describen las reglas que se recomiendan seguir para
seleccionar los valores de $a$, $X_0$ y $m$ dependiendo de si el sistema
en que se trabaja es binario o decimal.

\subsubsection*{a) Sistema decimal}

Si se trabaja en sistema decimal, los valores de los parámetros deben
ser seleccionados de acuerdo a los siguientes criterios:

\begin{enumerate}
\item El valor de la semilla puede ser cualquier entero impar no
divisible entre 2 ó 5 y debe ser relativamente primo a $m$.
\item El valor seleccionado de $a$ debe ser obtenido de acuerdo a la
siguiente identidad:
\begin{equation}
a = 200t \pm p
\end{equation}
donde $t$ es cualquier entero y $p$ es cualquiera de los siguientes
valores: 3, 11, 13, 19, 21, 27, 29, 37, 53, 59, 61, 67, 69, 77, 83, 91.
\item El valor seleccionado de $m$ puede ser $10^d$. Si $m=10^d$ y
$d \ge 5$ el período del generador es $5 \times 10^{d-2}$.
\end{enumerate}

Por otra parte, si $m=10^d$ y $d<5$, entonces el período del generador
se obtiene de acuerdo a la siguiente expresión:

\begin{equation}
\text{Período} = \text{Mínimo común múltiplo} \left\{ \lambda(P_1^{d_1}), \lambda(P_2^{d_2}), \ldots, \lambda(P_n^{d_n}) \right\}
\end{equation}

donde $P_i$ es un factor primo de $m$, y:

\begin{align*}
\lambda(2) = 1, \quad \lambda(4) = 2 \\
\lambda(2^d) = 2^{d-2} \text{ si } d \ge 3 \\
\lambda(p^d) = p^{d-1}(p-1) \text{ si } p \ge 2
\end{align*}

Con el propósito de ilustrar la obtención del período para este último
caso, analicemos el siguiente generador:

\[
X_{n+1} = 3X_n \bmod 100 \quad \text{y} \quad X_0 = 17
\]

puesto que $m$ puede ser expresado como $10^2$ o bien como $(2^2)(5^2)$,
entonces el período de este generador de acuerdo a la expresión anterior
sería:

\begin{align*}
\text{Período} &= \text{Mínimo común múltiplo} \{ (2^2), (5^2) \} \\
&= \text{Mínimo común múltiplo} \{2, 20\} \\
&= 20
\end{align*}

La tabla 2.4 muestra la secuencia de números pseudoaleatorios de este
generador. Como se puede apreciar en esta tabla, el período del
generador es 20.

\begin{table}[H]
\centering
\caption*{\textbf{Tabla 2.4} Números pseudoaleatorios del generador $X_{n+1} = 3X_n \bmod 100$}
\begin{tabular}{cccccccc}
\toprule
$n$ & $X_n$ & $n$ & $X_n$ & $n$ & $X_n$ & $n$ & $X_n$ \\
\midrule
1 & 51 & 6 & 93 & 11 & 99 & 16 & 57 \\
2 & 53 & 7 & 79 & 12 & 97 & 17 & 71 \\
3 & 59 & 8 & 37 & 13 & 91 & 18 & 13 \\
4 & 77 & 9 & 11 & 14 & 73 & 19 & 39 \\
5 & 31 & 10 & 33 & 15 & 19 & 20 & 17 \\
\bottomrule
\end{tabular}
\end{table}

\subsubsection*{b) Sistema binario}

Si se trabaja en sistema binario, los valores de los parámetros deben
ser seleccionados de acuerdo a los siguientes criterios:

\begin{enumerate}
\item El valor de la semilla puede ser cualquier entero impar
relativamente primo a $m$.
\item El valor seleccionado de $a$ debe ser obtenido a partir de la
siguiente expresión:
\begin{equation}
a = 8t \pm 3
\end{equation}
donde $t$ es cualquier entero.
\item El valor seleccionado de $m$ puede ser $2^d$. Si $m=2^d$ el
período del generador es $2^{d-2}$ ó $m/4$.
\end{enumerate}

Para ilustrar la obtención del período de un generador en sistema
binario, suponga que se tiene un generador en el cual los valores de sus
parámetros son: $a=5$, $X_0=5$ y $m=32$. Para estos valores, la
secuencia de números pseudoaleatorios son mostrados en la tabla 2.5.
Como se puede apreciar en esta tabla, el período del generador es 8.

\begin{table}[H]
\centering
\caption*{\textbf{Tabla 2.5} Números pseudoaleatorios del generador $X_{n+1} = 5X_n \bmod 32$}
\begin{tabular}{cccc}
\toprule
$n$ & $X_n$ & $n$ & $X_n$ \\
\midrule
1 & 25 & 5 & 9 \\
2 & 29 & 6 & 13 \\
3 & 17 & 7 & 1 \\
4 & 21 & 8 & 5 \\
\bottomrule
\end{tabular}
\end{table}

%==========================================================
\section*{PROBLEMAS}

\begin{enumerate}
\item[2.1.] Determine el período de los siguientes generadores
congruenciales mixtos:
\begin{align*}
X_{n+1} &= (8X_n + 16) \bmod 100 \text{ y } X_0 = 15 \\
X_{n+1} &= (50X_n + 17) \bmod 64 \text{ y } X_0 = 13 \\
X_{n+1} &= (5X_n + 24) \bmod 32 \text{ y } X_0 = 7 \\
X_{n+1} &= (5X_n + 21) \bmod 100 \text{ y } X_0 = 3 \\
X_{n+1} &= (9X_n + 13) \bmod 32 \text{ y } X_0 = 8
\end{align*}

\item[2.2.] Determine el período de los siguientes generadores
congruenciales multiplicativos:
\begin{align*}
X_{n+1} &= 203\,X_n \bmod 10^5 \text{ y } X_0 = 17 \\
X_{n+1} &= 211\,X_n \bmod 10^8 \text{ y } X_0 = 19 \\
X_{n+1} &= 221\,X_n \bmod 10^3 \text{ y } X_0 = 3 \\
X_{n+1} &= 5\,X_n \bmod 64 \text{ y } X_0 = 7 \\
X_{n+1} &= 11\,X_n \bmod 128 \text{ y } X_0 = 9
\end{align*}

\item[2.3.] Defina los parámetros $a$, $c$, $X_0$ y $m$ de un generador
congruencial mixto que aseguren período completo.

\item[2.4.] Defina los parámetros $a$, $X_0$ y $m$ de un generador
congruencial multiplicativo que aseguren período completo.

\item[2.5.] Defina la relación de recurrencia de los últimos $i$
dígitos de los siguientes generadores congruenciales mixtos:
\begin{align*}
X_{n+1} &= (21X_n + 221) \bmod 100,\ X_0 = 7 \text{ e } i = 1 \\
X_{n+1} &= (61X_n + 421) \bmod 1000,\ X_0 = 11 \text{ e } i = 2
\end{align*}
\end{enumerate}

\end{document}
